\section{Related Work}
\label{sec:related_work}

\subsection{Analytical and NWP-Based Altitude Harmonization}
Vertical navigation has historically operated on a dual-reference system. Manned aviation utilizes pressure altitude, which references a virtual datum (ISA MSL) to ensure relative separation between aircraft within the same air mass \cite{SIDs_2022}. In contrast, GNSS provides geometric altitude relative to the WGS-84 ellipsoid. The transformation between these systems involves correcting for two primary factors: the separation between the geoid and the ellipsoid (Geoid Undulation, $N$), and the deviation of the actual atmospheric temperature profile from the ISA model (Scale Factor Error) \cite{navi_637}. The ICARUS project highlighted that without real-time harmonization services, the mixing of barometric and geometric traffic in Very Low-Level (VLL) airspace could lead to catastrophic vertical separation infringements \cite{SIDs_2022}.

To bridge this gap, classic methods, such as the Blanchard algorithm~\cite{blanchard1971new}, integrate inertial measurements with air data to estimate errors induced by non-standard temperature and pressure gradients. However, these methods rely heavily on the accuracy of onboard temperature sensors and are prone to drift over long durations \cite{li2010errors}. More recently, researchers have leveraged Numerical Weather Prediction (NWP) data to act as a "virtual sensor". Simonetti and Garc\'{i}a Crespillo demonstrated that interpolating 4D weather data (e.g., ERA5) can reduce barometric altitude errors to meter-level~\cite{navi_637}. Similarly, Schumann \textit{et al.} utilized NWP analysis to verify aircraft height-keeping performance, proving that meteorological reanalysis can serve as a truth source for static pressure calibration \cite{Schumann_2017}. Despite their high accuracy, NWP-based methods are limited by the latency of data availability and the coarse spatial resolution of global models (typically $0.25^\circ$), which frequently fail to capture microclimate variations in urban canyons.

\subsection{Data-Driven Calibration for Barometric Sensors}
Conventional methods for barometric altitude estimation rely on the barometric formula with empirical corrections \cite{icao1993manual}. With the advent of deep learning, data-driven algorithms have been widely applied to model the non-linear relationship between pressure and geometric height. Recent work incorporates machine learning to learn sensor-specific biases, with ensemble approaches like Random Forests typically achieving 10 $\sim$ 30m accuracy in urban altitude estimation \cite{zhang2020urban}.

While standard neural networks and tree-based models can approximate these non-linear mappings, they act as ``black boxes'' and lack physical interpretability. Consequently, they often lead to poor generalization outside the training distribution (extrapolation), as they do not inherently respect the hydrostatic constraints governing atmospheric physics. Furthermore, these standard approaches traditionally lack explicit architectural mechanisms to decouple static unit hardware bias from the environmental model, inherently preventing zero-shot generalization across previously unseen, uncalibrated instruments in real-world deployments.

\subsection{Neural Fields and Spatial Representation}
To achieve continuous and robust spatial modeling, Coordinate-Based Neural Networks, or Neural Fields, have emerged as a powerful paradigm. Originally popularized by Neural Radiance Fields (NeRF) for novel view synthesis \cite{mildenhall2020nerf}, this architecture maps spatial coordinates to continuous physical quantities. To overcome the spectral bias of standard Multi-Layer Perceptrons (MLPs), M\"{u}ller \textit{et al.} introduced Instant-NGP, utilizing a multi-resolution hash encoding technique that dramatically accelerates the learning of high-frequency spatial details \cite{muller2022instant}.

While Neural Fields have revolutionized computer vision, their application to continuous geospatial regression and massive atmospheric parameter estimation remains largely underexplored. Furthermore, training Neural Fields on spatially heterogeneous multi-sensor IoT data presents severe convergence methodologies. Inspired by Bengio \textit{et al.} \cite{bengio2009curriculum}, Curriculum Learning—a strategy of training models on progressively harder examples—offers a potential solution. In this paper, we pioneer the adaptation of Multi-Resolution Hash Encoding, intertwined with SIREN features, for high-altitude metrological applications. By coupling an altitude-based curriculum learning scheme with explicit formulation of physical hardware pressure bias, we transform continuous Neural Fields into a highly accurate, zero-shot generalizable spatial measurement topology tool.
